\scanpage{10}
\mspage{9}

Les lieux d'une ficelle \inserted{orientée} forment un ens ordonné $\mathcal L$. Il satisfait les conditions suivantes
\margin{\unsure{Attention|Autre déf.} $X_{A,C}\subset X_{A,D}$ (et \deleted{$\unsure{C\le D}$} $\Leftrightarrow$ $\unsure{=}$ $X_{A,C}=X_{A,D}$) \unsure{n'implique} \uline{pas} que $C\le D$!}

a) $\exists$ plus petit et plus grand élément \inserted{$\alpha$ et $\beta$} distincts

b) $\forall x\in\mathcal L$, \deleted{\unsure{si} \unsure{$y,z$ sont des $y$}} \inserted{$x\ne\beta$, l'ens des élts $>x$ est} filtrant décroissant;

b') $\forall x\in\mathcal L$, $x\ne\alpha$, l'ens des \unsure{élts|éléments} $<x$ est filtrant croissant

c) Si $x<y$, $\exists z$ avec $x<z<y$ (divisibilité)

\uline{NB} \textcircled{1} Deux lieux \inserted{$x,y$} sont \uline{disjoints}, \deleted{ils} ssi $x<y$ ou $y<x$. Ils définissent une ficelle $F_{x,y}=F_{y,x}$.

\margin{\textcircled{2} Si on \unsure{prend l'or.} \unsure{opposée}, \deleted{\unsure{défini}} \unsure{ordre opposé}.}

\unsure{Considérons maintenant $n$} lieux intermédiaires, \uline{mutuellement disjoints} (i.e. \unsure{font|forment} une partie totalement ordonnée). La catégorie des syst. formés d'une ficelle orient. \inserted{$F$} et de $n$ tels lieux \inserted{$x_1\cdots x_n$} est équivalente à celle des systèmes $(F_1,\ldots,F_{n+1},h_{21},h_{32},\ldots,h_{n+1,n})$ de $n+1$ ficelles orientées, et de $n$ homeo. de raccordement.

\rotatebox{90}{\reflectbox{\unsure{EXTRA STRONG}}}
\rotatebox{90}{\reflectbox{\unsure{EXTRA STRONG}}}
\mspage{9}
